% GENERATED FILE. Edit canonical lessons, then run npm run generate:books.
% canonical-source-hash: fnv1a64:8061c08958cd531a
% canonical-lessons: ZH-C96-yinwei, ZH-W96-yin, ZH-C96-suoyi, ZH-C96-shufu, ZH-C96-practice

\chapter{Because}
\label{ch:zh-because}
\hlchaptermodality{96}

\begin{chapteropening}
\textbf{By the end of this chapter:} \emph{I can give one reason with \zh{因为}.}

Giving one reason with \zh{因为}, and what follows from it with suǒyǐ, built on one new written character: \zh{因}.
\end{chapteropening}

\section[\texorpdfstring{\zh{因为} (yīnwèi)}{因为 (yīnwèi)}]{\zh{因为} (yīnwèi) — because}

\label{lesson:ZH-C96-yinwei}



\begin{quote}
A new chapter, and it opens with the answer to the last one's question.

\begin{itemize}
\raggedright
  \item \emph{Recall:} say \zh{你为什么不去} — \emph{nǐ wèishénme bú qù}
  \item \emph{Recall:} say \emph{wǒ wàng le}, I forgot — \textbf{R1}, the chapter you have finished
  \item \emph{Recall:} say \emph{dǔchē}, a traffic jam — \textbf{R1}
\end{itemize}
\end{quote}

\subsection*{You'll want to know: \zh{因为}}
\textbf{\zh{因为}} — \emph{yīnwèi} — \textquotedblleft{}because\textquotedblright{}.

\begin{quote}
\textbf{\zh{我不去} \zh{因为我没有车}} \emph{wǒ bú qù, yīnwèi wǒ méiyǒu chē} — I'm not going, because I don't have a car.
\end{quote}

The next lesson puts the first character in your hand.

\begin{cousinweb}
\textbf{\zh{因为}} (\emph{yīnwèi}) is built from \zh{因} 'to follow from' and \zh{为} 'for'. \zh{因} is inherited from Middle Chinese \emph{'jin} (Guangyun, in Baxter's transcription), glossed 'to rely on; to follow from'.
\end{cousinweb}

\begin{grammarlens}[title={What you've built}]
One reason, in one clause, after what it explains: \zh{我不去}, \zh{因为我没有车}. It can also come first: \zh{因为我没有车}, \zh{我不去}.
\end{grammarlens}

\subsection*{Your turn}
\begin{itemize}
\raggedright
  \item \emph{Say it:} \emph{yīnwèi}
  \item \emph{Say it:} \emph{yīnwèi}, once more
  \item \emph{Say it:} \zh{我不去} \zh{因为我没有车} — \emph{wǒ bú qù, yīnwèi wǒ méiyǒu chē} — I'm not going, because I don't have a car.
  \item \emph{Say it:} \zh{因为我没有车} \zh{我不去} — \emph{yīnwèi wǒ méiyǒu chē, wǒ bú qù} — Because I don't have a car, I'm not going.
  \item \emph{Recall:} say \emph{nǐ wèishénme bú qù}, then answer it with \zh{因为}
\end{itemize}

\begin{tcolorbox}[breakable,colback=teal!4,colframe=teal!35!black,title={Before you move on}]
What does \zh{因为} mean? (\textquotedblleft{}Because\textquotedblright{}.) Which question does it answer? (\textbf{\zh{为什么}}.)
\end{tcolorbox}
\section[\texorpdfstring{\zh{因} (yīn)}{因 (yīn)}]{\zh{因} — \zh{大} inside the frame of \zh{国}}

\label{lesson:ZH-W96-yin}



\begin{quote}
The word first, then the character inside it.

\begin{itemize}
\raggedright
  \item \emph{Recall:} say \emph{yīnwèi} — \textbf{R1}, one lesson back, before it decays
\end{itemize}

\begin{quote}
\zh{国} \zh{大}
\end{quote}

\emph{Read it:} each one aloud, with nothing said first — \emph{guó}, \emph{dà}
\end{quote}

\subsection*{Script — \zh{大} inside the frame of \zh{国}}
\hlblockfigure{\detokenize{figures/ZH-W96-yin-filmstrip.pdf}}{How \zh{因} is written, stroke by stroke}

\begin{quote}
\zh{因}
\end{quote}

\textbf{A frame and what it holds, and the frame is closed last, after its contents, as in \zh{国}.}

\textbf{The frame, two strokes first}: its left side, then its top turning down the right side.

\textbf{\zh{大}, three strokes}, inside: the bar, the left-falling stroke, the short right-falling stroke.

\textbf{The bottom bar}, last, closing the frame.

2 + 3 + 1 = six.

\subsection*{Your turn}
\begin{itemize}
\raggedright
  \item \emph{Write it:} \zh{因}, all six, with the model in front of you
  \item \emph{Cover:} write it again from memory, and count the pieces, not the strokes
\end{itemize}

\begin{tcolorbox}[breakable,colback=teal!4,colframe=teal!35!black,title={Before you move on}]
How many strokes in \zh{因}? (\textbf{Six}.) What is inside the frame? (\textbf{\zh{大}}.) When is the frame closed? (\textbf{Last, as in \zh{国}}.)

Source: \href{https://raw.githubusercontent.com/chanind/hanzi-writer-data/68d10a4b21150cae5e1ebbd223eed289cf32d90c/data/\%E5\%9B\%A0.json}{Hanzi Writer Data \zh{因}, pinned stroke paths and medians}.
\end{tcolorbox}
\section[suǒyǐ]{suǒyǐ — so; therefore}

\label{lesson:ZH-C96-suoyi}



\begin{quote}
Before the new one: read the last lesson's character, then two words, one of which ends like this one.

\begin{itemize}
\raggedright
  \item \emph{Read it:} \zh{因为}, aloud — \emph{yīnwèi}
  \item \emph{Recall:} say \emph{lèi}, tired — \textbf{R4}, recognition at distance
  \item \emph{Recall:} say \emph{kěyǐ}, may — \textbf{R2}
\end{itemize}
\end{quote}

\subsection*{You'll want to know: suǒyǐ}
\emph{suǒyǐ} — \textquotedblleft{}so; therefore\textquotedblright{}.

\begin{quote}
\emph{yīnwèi wǒ hěn lèi, suǒyǐ wǒ bú qù} — I'm tired, so I'm not going.
\end{quote}

Two dips in a row: spoken, the first rises, \emph{suóyǐ}. It is written \emph{suǒyǐ}.

Its characters come later; for now it is a word for the ear.

\begin{cousinweb}
\emph{suǒyǐ} is built from suǒ 'a place; that which' and yǐ 'to use', the yǐ of \emph{kěyǐ}. Both are inherited from Middle Chinese, \emph{srjoX} and \emph{yiX} (Guangyun, in Baxter's transcription), \emph{srjoX} glossed 'a place'. The regular change from \emph{srjoX} would give shǔ, not suǒ, so why it is said suǒ is not settled.
\end{cousinweb}

\begin{grammarlens}[title={What you've built}]
\zh{因为} gives the reason, \emph{suǒyǐ} what follows from it. Mandarin often says both: \emph{yīnwèi …, suǒyǐ …}. Either half can also stand alone.
\end{grammarlens}

\subsection*{Your turn}
\begin{itemize}
\raggedright
  \item \emph{Say it:} \emph{suǒyǐ} — and hear it as \emph{suóyǐ}
  \item \emph{Say it:} \emph{yīnwèi wǒ hěn lèi, suǒyǐ wǒ bú qù} — I'm tired, so I'm not going.
  \item \emph{Recall:} say \emph{kěyǐ}, then \emph{suǒyǐ}, and hear the same yǐ
\end{itemize}

\begin{tcolorbox}[breakable,colback=teal!4,colframe=teal!35!black,title={Before you move on}]
What does \emph{suǒyǐ} mean? (\textquotedblleft{}So; therefore\textquotedblright{}.) Which half gives the reason, \zh{因为} or \emph{suǒyǐ}? (\textbf{\zh{因为}}.)
\end{tcolorbox}
\section[shūfu]{shūfu — comfortable; well}

\label{lesson:ZH-C96-shufu}



\begin{quote}
Before the new one: last lesson's word, then two words about health.

\begin{itemize}
\raggedright
  \item \emph{Recall:} say \emph{suǒyǐ}, so — \textbf{R1}, one lesson back
  \item \emph{Recall:} say \emph{shēntǐ}, health — \textbf{R4}, recognition at distance
  \item \emph{Recall:} say \emph{bìng}, to fall ill — \textbf{R2}
\end{itemize}
\end{quote}

\subsection*{You'll want to know: shūfu}
\emph{shūfu} — \textquotedblleft{}comfortable; well\textquotedblright{}.

\begin{quote}
\emph{wǒ bù shūfu} — I'm not feeling well.
\end{quote}

\begin{quote}
\emph{zhège yǐzi hěn shūfu} — This chair is comfortable.
\end{quote}

Its characters come later; for now it is a word for the ear.

\begin{cousinweb}
\emph{shūfu} is built from shū 'to stretch; to ease' and fu, the fu of \emph{yīfu}. Both are inherited from Middle Chinese, \emph{syo} and \emph{bjuwk} (Guangyun, in Baxter's transcription), \emph{syo} glossed 'slow; to stretch out'.
\end{cousinweb}

\begin{grammarlens}[title={What you've built}]
\emph{wǒ bù shūfu} is the gentle way to say you are unwell, and a reason nobody argues with: \emph{yīnwèi wǒ bù shūfu}.
\end{grammarlens}

\subsection*{Your turn}
\begin{itemize}
\raggedright
  \item \emph{Say it:} \emph{shūfu}
  \item \emph{Say it:} \emph{shūfu}, once more
  \item \emph{Say it:} \emph{wǒ bù shūfu} — I'm not feeling well.
  \item \emph{Say it:} \emph{zhège yǐzi hěn shūfu} — This chair is comfortable.
  \item \emph{Recall:} say \emph{bìng}, then \emph{shūfu} again
\end{itemize}

\begin{tcolorbox}[breakable,colback=teal!4,colframe=teal!35!black,title={Before you move on}]
What does \emph{shūfu} mean? (\textquotedblleft{}Comfortable; well\textquotedblright{}.) Say you are not feeling well. (\emph{wǒ bù shūfu}.)
\end{tcolorbox}
\section[Practice]{\zh{为什么} — \zh{因为}}

\label{lesson:ZH-C96-practice}



\begin{quote}
The chapter closes here, with every why answered.

\begin{itemize}
\raggedright
  \item \emph{Recall:} say \emph{shūfu}, comfortable — \textbf{R1}, one lesson back
  \item \emph{Recall:} say \emph{è}, hungry — \textbf{R4}, recognition at distance
  \item \emph{Recall:} say \zh{上课} — \emph{shàngkè}, to go to class — \textbf{R2}
\end{itemize}
\end{quote}

\subsection*{The exchange}
\begin{quote}
\textbf{A:} \zh{你为什么不吃饭}
\end{quote}

\begin{quote}
\emph{nǐ wèishénme bù chī fàn} — Why aren't you eating?
\end{quote}

\begin{quote}
\textbf{B:} \emph{yīnwèi wǒ bú è} — Because I'm not hungry.
\end{quote}

\begin{quote}
\textbf{A:} \zh{你为什么不去上课}
\end{quote}

\begin{quote}
\emph{nǐ wèishénme bú qù shàngkè} — Why aren't you going to class?
\end{quote}

\begin{quote}
\textbf{B:} \emph{yīnwèi wǒ bù shūfu, suǒyǐ wǒ bú qù} — I'm not feeling well, so I'm not going.
\end{quote}

\begin{quote}
\textbf{A:} \zh{你明天能来吗}
\end{quote}

\begin{quote}
\emph{nǐ míngtiān néng lái ma} — Can you come tomorrow?
\end{quote}

\begin{quote}
\textbf{B:} \zh{能}
\end{quote}

\begin{quote}
\emph{néng} — I can.
\end{quote}

\zh{为什么} asks; \zh{因为} answers, with one reason in one clause; \emph{suǒyǐ} carries the reason on to what happens.

\begin{tcolorbox}[breakable,skin=enhanced,colback=orange!6,colframe=orange!45!black,boxrule=0.5pt,arc=1mm,left=6pt,right=6pt,top=4pt,bottom=4pt,fonttitle=\bfseries,title={Read this}]
\emph{Read it:} the exchange through once without stopping, then answer the question below
\end{tcolorbox}

\begin{grammarlens}[title={What you've built — a question and its reason}]
\begin{itemize}
\raggedright
  \item \textbf{\zh{你为什么不去}} \emph{nǐ wèishénme bú qù} — why aren't you going?
  \item \textbf{\zh{因为我没有车}} \emph{yīnwèi wǒ méiyǒu chē} — because I don't have a car.
  \item \emph{yīnwèi wǒ hěn lèi, suǒyǐ wǒ bú qù} — I'm tired, so I'm not going.
\end{itemize}

\textbf{Listen.} \emph{Nǐ wèishénme xuéxí Hànyǔ?}

\textbf{Speak.} Answer it with \zh{因为} and one reason of your own.

\emph{Write it:} \zh{因} — the frame's two sides, \zh{大} inside, then the bar that closes it
\end{grammarlens}

\subsection*{Your turn}
\begin{itemize}
\raggedright
  \item \emph{Say it:} \zh{你为什么不去} — \emph{nǐ wèishénme bú qù} — Why aren't you going?
  \item \emph{Say it:} \zh{因为我没有车} — \emph{yīnwèi wǒ méiyǒu chē} — Because I don't have a car.
  \item \emph{Say it:} the whole exchange, both parts
  \item \emph{Write it:} \zh{因为}
\end{itemize}

\begin{tcolorbox}[breakable,colback=teal!4,colframe=teal!35!black,title={Before you move on}]
Give a reason with no car in it. (\textbf{\zh{因为我没有车}}.) Say you are not well, so you are not going. (\emph{yīnwèi wǒ bù shūfu, suǒyǐ wǒ bú qù}.) Which question does \zh{因为} answer? (\textbf{\zh{为什么}}.)
\end{tcolorbox}
